\chapter{Теория}

\section{Разделение на сети по 2000, 1000, 500, 200, 100, 50 узла}


\begin{figure}[H] 
    \centering
    \includegraphics[width=0.75\textwidth]{1.png}
    \caption{Начальные данные}
\end{figure}

\subsection{Выделение сети на 2000 узлов}
Чтобы сеть могла вместить 2000 узлов, нужно оставить в узловой части 11 бит, тогда: \\

Адрес сети: 10.1.0.0 \\
Диапазон адресов: 10.1.0.1-10.1.7.254

\subsection{Выделение сети на 1000 узлов}
Чтобы сеть могла вместить 1000 узлов, нужно оставить в узловой части 10 бит, тогда: \\

Адрес сети: 10.1.8.0 \\
Диапазон адресов: 10.1.8.1-10.1.11.254

\subsection{Выделение сети на 500 узлов}
Чтобы сеть могла вместить 500 узлов, нужно оставить в узловой части 9 бит, тогда: \\

Адрес сети: 10.1.12.0 \\
Диапазон адресов: 10.1.12.1-10.1.13.254

\subsection{Выделение сети на 200 узлов}
Чтобы сеть могла вместить 200 узлов, нужно оставить в узловой части 8 бит, тогда: \\

Адрес сети: 10.1.14.0 \\
Диапазон адресов: 10.1.14.1-10.1.14.254

\subsection{Выделение сети на 100 узлов}
Чтобы сеть могла вместить 100 узлов, нужно оставить в узловой части 7 бит, тогда: \\

Адрес сети: 10.1.15.0 \\
Диапазон адресов: 10.1.15.1-10.1.15.126

\subsection{Выделение сети на 50 узлов}
Чтобы сеть могла вместить 50 узлов, нужно оставить в узловой части 6 бит, тогда: \\

Адрес сети: 10.1.15.128 \\
Диапазон адресов: 10.1.15.129-10.1.15.190
\endinput